\section{Ejercicios 10}

\begin{ejercicio}
\[
3^n+5^n>2^{n+2}
\]

Probar por inducción la proposición para todo $n\geq2$.

\textbf{Caso base:} $n=2$.

\[
3^2+5^2\geq2^{2+2}
\]

\[
9+25\geq16
\]

\[
34\geq16
\]

Por lo tanto, el caso base es verdadero.

\textbf{Paso inductivo:}

Supongamos que la proposición es verdadera para $n=k$, es decir:

\textbf{Hipótesis inductiva:}

\[
3^k+5^k>2^{k+2}.
\]

Queremos demostrar que la proposición es verdadera para $n=k+1$.

\textbf{Tesis inductiva:}

\[
3^{k+1}+5^{k+1}>2^{(k+1)+2}.
\]

Veamos:

\[
3^{k+1}+5^{k+1}
=
3\cdot3^k+5\cdot5^k.
\]

Como $3>2$ y $5>2$, tenemos:

\[
3\cdot3^k>2\cdot3^k
\]

y

\[
5\cdot5^k>2\cdot5^k.
\]

Por lo tanto,

\[
3\cdot3^k+5\cdot5^k
>
2\cdot3^k+2\cdot5^k.
\]

Entonces,

\[
3^{k+1}+5^{k+1}
>
2(3^k+5^k).
\]

Por la hipótesis inductiva,

\[
3^k+5^k>2^{k+2}.
\]

Como $2>0$, podemos multiplicar por $2$:

\[
2(3^k+5^k)>2\cdot2^{k+2}.
\]

Por lo tanto,

\[
3^{k+1}+5^{k+1}
>
2\cdot2^{k+2}.
\]

Finalmente,

\[
2\cdot2^{k+2}=2^{k+3}=2^{(k+1)+2}.
\]

Por lo tanto,

\[
3^{k+1}+5^{k+1}>2^{(k+1)+2}.
\]

Así, la tesis inductiva es verdadera.

Como el caso base es verdadero y hemos demostrado que si la proposición es verdadera para $n=k$, entonces también lo es para $n=k+1$, queda demostrado por inducción que

\[
\boxed{3^n+5^n>2^{n+2}}
\]

para todo $n\geq2$.
\end{ejercicio}

\begin{ejercicio}
\[
p(n): 3^n \geq n^3 
\]

Entonces pruebo por inducción que $p(n)$ es verdadero para todo $n\geq 1$.

\textbf{Caso base:} $p(1)$

\[
3^1 \geq 1^3
\]

\[
3 \geq 1 
\]

verdadero.

\textbf{Paso Inductivo:} $p(k)$

\vspace{0.5cm}
supongo que $p(k)$ es verdadero, es decir:

\textbf{Hipótesis inductiva:} 

\[
3^k \geq k^3
\]

Quiero probar que $p(k+1)$ es verdadero:

\textbf{Tesis inductiva:}

\[
3^{k+1} \geq (k+1)^3
\]

Veamos:

\[ 
3^{k+1} = 3 \cdot 3^k
\]
es igual termino de la hipótesis inductiva:
\[ 
HI \geq 3 \cdot k^3 \geq  (k+1)^3
\]

\[ 
\geq 3 \cdot k^3 \geq k^3 + 3k^2 + 3k + 1
\]

\[ 
\geq 3k^3 \geq k^3 + 3k^2 + 3k + 1
\]

\[ 
\geq 2k^2 - 3k^1 - 3k \geq 1
\]

\[ 
\geq k(k^2 - 3k - 3) \geq 1
\]

\begin{solucion}{raices de la ecuación cuadrática}
aca se puede ver que la ecuación cuadrática $k^2 - 3k - 3 = 0$ tiene dos raices, una negativa y otra positiva, por lo tanto para todo $k \geq 4$ se cumple que $k^2 - 3k - 3 \geq 0$, entonces para todo $k \geq 4$ se cumple que $k(k^2 - 3k - 3) \geq 1$.
solo nos queda verificar los casos $k=1,2,3$ de los cuales ya sabemos que $p(1)$ es verdadero, y verificando $p(2)$ y $p(3)$ se cumple que son verdaderos, por lo tanto queda demostrado que $p(n)$ es verdadero para todo $n \geq 1$.
\end{solucion}



\end{ejercicio}